\section{展開と割り算}\label{節:展開と割算}

$x(2x-1)^2$ や $\frac{x^3+2x^2+3}{x+1}$ は,そのままでは前節までの表のどれにも当てはまらない.しかし,式を同じ値の別の式に書き換えれば,\,$x$ のべきや一次式の逆数の和にできる.本節では,展開と割り算という2つの書き換えによって,基本公式に帰着させる.

ここで1つ断りをいれておこう.割り算や約分で式を書き換えても,積分するのはもとの関数である.たとえば $\frac{x^2-1}{x-1}=x+1$ は $x\ne1$ でしか成り立たない.右辺が $x=1$ でも定義されるからといって,左辺の関数がその点で定義されたことにはならない.したがって,書き換えた後に得た答えも,もとの関数の分母が $0$ にならない各区間で読む.\footnote{積分する区間にその点が含まれる場合は,連続に延長した関数を扱うのか,もとの式を扱うのかを明らかにする必要がある.}

まず整式の積である.積の形では微分の逆が見えなくても,展開すれば各項が $x$ のべきになり,基本公式を1項ずつ使える.式を短く保つことと,積分しやすくすることは同じではないのだ.

次に,整式 $P(x),Q(x)$ の商 $\frac{P(x)}{Q(x)}$ を考える.分子の次数が分母以上のときは,整式の割り算をして
\[
 P(x)=Q(x)S(x)+R(x)\qquad(R(x)\text{ の次数}<Q(x)\text{ の次数})
\]
と書く.両辺を $Q(x)$ で割ると
\[
 \frac{P(x)}{Q(x)}=S(x)+\frac{R(x)}{Q(x)}
\]
となり,分数全体が「整式」と「分子の次数が低い分数」の和に分かれる.整式 $S(x)$ はべきの公式で積分できる.特に分母が一次式 $x-\alpha$ のときは,余りは定数 $r$ であり,\footnote{剰余の定理より $r=P(\alpha)$ である.割り算の検算に使える.}
\[
 \int\frac{P(x)}{x-\alpha}\dd x=\int S(x)\dd x+r\log|x-\alpha|+C
\]
となる.

\begin{bookpoint}
\textbf{POINT}\par
整式の積は展開し,分子の次数が分母以上の分数は $P=QS+R$ と割って,整式 $S$ と $\frac RQ$ の和に分ける.
\end{bookpoint}

\begin{example}{積を展開する理由}
$\int x(2x-1)^2\dd x$ を求めよ.
\end{example}

\paragraphbf{方針}
2乗を展開して $x$ を掛け,\,$x$ のべきの和にする.

\paragraphbf{解答}
$x^4-\frac43x^3+\frac12x^2+C$

\paragraphbf{解法}
2乗を展開してから $x$ を掛けると
\begin{align*}
 x(2x-1)^2
 &=x(4x^2-4x+1)\\
 &=4x^3-4x^2+x
\end{align*}
となる.したがって
\[
 \int x(2x-1)^2\dd x
 =x^4-\frac43x^3+\frac12x^2+C
\]
である.整式なので,実数全体で成り立つ.右辺を微分すれば $4x^3-4x^2+x$ に戻り,これをくくり直すと元の式になる.

\begin{example}{分子の次数を下げる}
$x\ne-1$ において,\,$\int\frac{x^3+2x^2+3}{x+1}\dd x$ を求めよ.
\end{example}

\paragraphbf{方針}
分子を分母 $x+1$ で割り,整式と $\frac{r}{x+1}$ の和に分ける.

\paragraphbf{解答}
$\frac{x^3}{3}+\frac{x^2}{2}-x+4\log|x+1|+C$

\paragraphbf{解法}
分子を分母で割ると
\[
 x^3+2x^2+3=(x+1)(x^2+x-1)+4
\]
となる.余りは,分子に $x=-1$ を代入した $-1+2+3=4$ と一致する.両辺を $x+1$ で割って
\[
 \frac{x^3+2x^2+3}{x+1}
 =x^2+x-1+\frac4{x+1}
\]
とすれば,すべて既知の形である.よって
\[
 \int\frac{x^3+2x^2+3}{x+1}\dd x
 =\frac{x^3}{3}+\frac{x^2}{2}-x+4\log|x+1|+C
\]
を得る.もとの関数は $x=-1$ で定義されないので,この答えは $x<-1$ と $x>-1$ の各区間で成り立つ.

\paragraphbf{展開が向かない場合}
展開は,数項で済むならば見通しのよい方法である.しかし,高いべきでは展開する項数が多くなるため,別の方法が有利になる.これは本章の終わりで改めて取り上げる.また,割り算を済ませても,余り $\frac{R(x)}{Q(x)}$ の分母が2次以上のことがある.その分数を分けるのが次節の部分分数である.
